\chapter{Envelope Generator}
\label{ch:envelope}

Each voice has an envelope generator that shapes the loudness of its sound over time. The envelope has a level from 0
to 4095, a 12-bit unsigned number, and four phases: attack, decay, sustain and release. The four values that shape it
are in the instrument record, so every instrument has its own shape (Table~\ref{tab:instrument}).

\begin{figure}[H]
	\centering
	\begin{tikzpicture}[
			scale=1.0,
			axis/.style={-{Stealth[length=2mm]}, thick},
			note/.style={font=\scriptsize, text=sircgray}
		]
		\draw[axis] (0,0) -- (9.4,0) node[right, font=\scriptsize] {time};
		\draw[axis] (0,0) -- (0,3.7) node[above, font=\scriptsize] {level};
		\draw[very thick] (0,0) -- (1.4,3.2) -- (2.9,1.8) -- (6.0,1.8) -- (8.4,0);
		\draw[dashed, sircgray] (1.4,3.2) -- (1.4,0);
		\draw[dashed, sircgray] (2.9,1.8) -- (2.9,0);
		\draw[dashed, sircgray] (6.0,1.8) -- (6.0,0);
		\draw[dashed, sircgray] (0,3.2) -- (1.4,3.2);
		\draw[dashed, sircgray] (0,1.8) -- (2.9,1.8);
		\node[note, anchor=east] at (0,3.2) {4095};
		\node[note, anchor=east] at (0,1.8) {sustain};
		\node[note] at (0.7,-0.3) {attack};
		\node[note] at (2.15,-0.3) {decay};
		\node[note] at (4.45,-0.3) {sustain};
		\node[note] at (7.2,-0.3) {release};
		\node[note, anchor=south] at (6.0,0.05) {stop};
	\end{tikzpicture}
	\caption{The Envelope Phases}
	\label{fig:envelope}
\end{figure}

\section{Phases}
\label{sec:envelope-phases}

The envelope is updated once in every 64 output samples (an interval of about 2.0 ms, or 499.4 updates a second). At
each update it takes one step, as follows:

\begin{table}[H]
	\centering
	\begin{tabular}{lp{9cm}}
		\toprule
		\textbf{Phase} & \textbf{Behaviour at each update}                                                                                                                                             \\
		\midrule
		Attack         & The level rises by \texttt{ATTACK}. When it reaches 4095, it stops there and the phase becomes decay. \texttt{ATTACK} of 0 sets the level to 4095 at once                     \\
		Decay          & The level falls by \texttt{DECAY}, but not below the sustain level. When it reaches that level, the phase becomes sustain. \texttt{DECAY} of 0 sets the sustain level at once \\
		Sustain        & The level stays at the sustain level until the voice is stopped                                                                                                               \\
		Release        & The level falls by \texttt{RELEASE}. When it reaches 0 the envelope is inactive. \texttt{RELEASE} of 0 sets the level to 0 at once                                            \\
		Inactive       & The level is 0 and stays there. The voice is silent until a note starts                                                                                                       \\
		\bottomrule
	\end{tabular}
	\caption{Envelope Phases}
	\label{tab:envelope-phases}
\end{table}

\texttt{ATTACK}, \texttt{DECAY} and \texttt{RELEASE} are unsigned 8-bit steps in units of the 12-bit level. \texttt{SUSTAIN}
is an unsigned 8-bit level, converted to the 12-bit scale by repeating its top four bits below it: the 12-bit level is
$(\texttt{SUSTAIN} \ll 4) \mathbin{|} (\texttt{SUSTAIN} \gg 4)$. A \texttt{SUSTAIN} of 255 gives the full level, 4095, and 0
gives 0. When the sustain level is 0 the voice is silent once the decay ends, but it is still active as far as the rules of Section~\ref{sec:note-start} go: its envelope is in the sustain phase, not inactive.

A stop (Section~\ref{sec:reg-command}) puts the envelope in the release phase from whatever phase it is in. A stop on
an inactive envelope has no effect.

\section{Starting a Note}

A note that retriggers (Section~\ref{sec:note-start}) sets the level to 0 and the phase to attack, and takes the first
step at once, without waiting for the voice's next update. An \texttt{ATTACK} of 0 therefore reaches full level
immediately, and any other attack starts at its first step and not at 0. The voice's own update then continues at its
regular intervals.

A full-scale attack takes at most $(\lceil 4095 / \mathrm{step} \rceil - 1) \times 2.0$ ms after the note starts, since
the first regular update falls within one interval of the note. A full-scale decay or release takes at most $\lceil
	4095 / \mathrm{step} \rceil \times 2.0$ ms. Table~\ref{tab:envelope-times} gives $4095 / \mathrm{step}$ intervals of
2.0 ms.

\begin{table}[H]
	\centering
	\begin{tabular}{rr}
		\toprule
		\textbf{Step} & \textbf{Time for 4095 levels (about)} \\
		\midrule
		255           & 32 ms                                 \\
		128           & 64 ms                                 \\
		64            & 128 ms                                \\
		32            & 256 ms                                \\
		16            & 512 ms                                \\
		8             & 1.0 s                                 \\
		4             & 2.05 s                                \\
		2             & 4.1 s                                 \\
		1             & 8.2 s                                 \\
		\bottomrule
	\end{tabular}
	\caption{Envelope Step and Time}
	\label{tab:envelope-times}
\end{table}

A step of 1 is the slowest the envelope goes. A voice with an \texttt{ATTACK} of 1 is silent at first: with a step of 1
the level is 1, which is too small to produce any output after the volume is applied (Section~\ref{sec:voice-gain}),
and the voice becomes audible only after one or more updates, later for a lower volume.

\section{Update Timing}
\label{sec:envelope-timing}

The six voices do not update together. Their updates are spread evenly across the 64 output samples of each interval,
so that in any one output sample at most one voice updates, and the arithmetic of the updates does not pile up. Each
voice still updates once in 64 output samples; the stagger fixes only which sample of the interval it is. The
assignment is fixed and is the same after every reset: voice $n$ updates when the count of sample periods since reset,
modulo 64, equals $(64 - \lfloor 64n/6 \rfloor) \bmod 64$, which is 0, 54, 43, 32, 22 and 11 for voices 0 to 5.

An envelope update is visible to the mixer from the next output sample.

\section{Retriggering}

A note that retriggers a voice that is still active restarts its envelope at 0 (Section~\ref{sec:note-start}). The
sound of the old note is replaced at once by the start of the new one, with no crossfade. A new note, with a position
of 0, begins with the first byte of its sample at its new level, so the output moves in one sample from wherever the
old note was to the new note's first value. A tune that retriggers a voice while the previous note is still loud has a
step in the output at each note. Chapter~\ref{ch:programming} describes how an instrument is designed to avoid it.
